\documentclass[10pt,a4paper]{amsart}
\usepackage[margin=19mm]{geometry}
\usepackage{amsmath,amssymb,mathtools}
\usepackage[scr=rsfs,cal=euler]{mathalfa}
\usepackage{xfrac}
\usepackage[unicode,hidelinks,bookmarks=false]{hyperref}
\newcommand{\ie}{i.e.\ }
\newcommand{\cf}{cf.\ }
\newcommand{\resp}{respectively\ }
\newcommand{\Subsection}{\subsection}
\providecommand{\nobreakdash}{\nobreak\hbox{-}}
\setlength{\emergencystretch}{2em}
\multlinegap0pt
\usepackage{bm}
\usepackage{bbm}
\usepackage{dsfont}
\usepackage{xspace}
\usepackage{cancel}
\usepackage{mathtools}
\usepackage{amsthm}
\makeatletter
\@ifundefined{thm}{\newtheorem{thm}{Theorem}}{}
\makeatother
\title[On Nieuwland Numbers and Polar Duality]{On Nieuwland Numbers and Polar Duality}
\author{Kavin Satheeskumar, Liam Benoit}
\date{Source: arXiv:2608.14912v1 (CC BY 4.0)}
\begin{document}
\maketitle

\noindent\emph{Source and licence.} On Nieuwland Numbers and Polar Duality. Version arXiv:2608.14912v1 (CC BY 4.0), distributed under the Creative Commons Attribution 4.0 International licence, as recorded in the arXiv licence metadata for this exact version and shown on the abstract page https://arxiv.org/abs/2608.14912v1. Full rights notice: licenses/arxiv-2608.14912-CC-BY-4.0.txt.\par\smallskip

\noindent\textit{Editorial scope.} Counted unit: the Thm whose source label is \texttt{thm:pass\_through\_simplified} in the article (source0.tex lines 419--432, source label \texttt{thm:pass\_through\_simplified}), delivered together with the article's own complete original proof of it (source0.tex lines 433--523, one \texttt{proof} environment of 91 lines). No mathematics is written, repaired or re-derived: statement and proof are the article's own text line for line. Required context retained uncounted: none. Every definition, display and label of those lines is retained unchanged. Omitted from this package and not redistributed: 1--418, 524--747. Nothing inside a retained excerpt is omitted or altered: every retained body is delivered exactly as the article's lines are written, so no omission receipt of any kind accompanies this package. Citations: the retained excerpts carry no citation command at all, so no bibliography entry of the article is transcribed and no reference is redistributed. Reference adaptation: none was needed and none was made; every reference inside the delivered unit resolves inside it, so no unresolved reference or citation is produced. Printed slips kept literal: no printed word, symbol, hypothesis or number of the counted statement or of its proof was altered. Layout and class: the article's own class is \texttt{amsart} with packages including inputenc, amsthm, amsmath, amssymb, mathrsfs, geometry, comment, graphicx, mathtools, tikz, tikz-3dplot, algorithm, algpseudocode, realboxes; the wrapper is the lane's pinned \texttt{amsart} editorial wrapper at 10pt on A4, which restates the theorem-like environments the retained text needs and adds the article's own macro definitions that the retained text needs (none), with extra wrapper packages (bm, bbm, dsfont, xspace, cancel, mathtools). The delivered numbering is the unit's own. Assets: no retained span contains a figure environment, a graphics-inclusion command, a PDF-page inclusion, a table or a TikZ picture; no image file is copied into this package, the package declares no asset dependency and no image file is redistributed with it. Licence: version arXiv:2608.14912v1 of this article is distributed under the Creative Commons Attribution 4.0 International licence, as recorded in the arXiv licence metadata for this exact version and shown on the abstract page https://arxiv.org/abs/2608.14912v1. A full rights notice is delivered at licenses/arxiv-2608.14912-CC-BY-4.0.txt. Characters: every retained excerpt is pure ASCII, so the delivered engine, XeLaTeX with its default Latin Modern text font, reports no missing character.
\begin{thm} \label{thm:pass_through_simplified}
	Let $P=\textup{conv}(v_1 \cdots v_n)$ and $Q=\{Ax \leq b\}$ be polytopes and let the rows of $A$ be $a_1^T \cdots a_m^T$. Let $\eta \neq 0$, then the following are equivalent.
	\begin{itemize}
		\item $P$ passes through $Q$ with certificate $(U, \eta, \delta)$
		\item There exists a bipartition $(I_+, I_-)$ of the set $\{1 \cdots m\}$ such that the following are true.
		      \begin{enumerate}
		      	\item For all $i \in I_+, a_i^T \eta \geq 0$ and for all $j \in I_-, a_j^T \eta \leq 0$.
		      	\item For all $1 \leq k \leq n, i \in I_+, j \in I_-$
		      	      \begin{gather*}
		      	      	(b_j-a_j^T(Uv_k+\delta))(a_i^T\eta) \geq (b_i-a_i^T(Uv_k+\delta))(a_j^T\eta).
		      	      \end{gather*}
		      \end{enumerate}
	\end{itemize}
\end{thm}

\begin{proof}
	$(\Longrightarrow)$
	\\\\
	By the assumption there are $z_1, \cdots, z_n$ such that for all $1 \leq k \leq n$, we have
	\begin{gather*}
		A(Uv_k + \delta+z_k\eta) \leq b,\\
		A(Uv_k + \delta)+z_kA\eta \leq b,\\
		\begin{bmatrix}A & A\eta\end{bmatrix}\begin{bmatrix}Uv_k + \delta \\ z_k\end{bmatrix} \leq b.
	\end{gather*}
	Let $I_- = \{1 \leq i \leq m:a_i^T\eta < 0\}$ and $I_+ = \{1 \leq i \leq m:a_i^T\eta \geq 0\}$. Then $(I_-, I_+)$ is a bipartition of $\{1 \cdots m\}$.\\\\
	Now consider arbitrary $i, j, k$ such that $1 \leq k \leq n$, $i \in I_+$ and $j \in I_-$. We prove that $(I_-, I_+)$ satisfies the conclusions of the theorem. We consider two cases, $a_i^T\eta = 0$ and $a_i^T\eta > 0$. Recall that $a_j^T\eta < 0$ by the construction of $I_+$ and $I_-$.\\\\
	Consider the case where $a_i^T\eta = 0$. By the assumption, we know
	\begin{gather*}
		a_i^T(Uv_k+\delta+z_k \eta) \leq b_i.
	\end{gather*}
	Since $a_i^T \eta = 0$, this becomes
	\begin{gather*}
		a_i^T(Uv_k+\delta+z_k \eta) \leq b_i,\\
		a_i^T(Uv_k+\delta)+z_k a_i^T\eta \leq b_i,\\
		a_i^T(Uv_k+\delta) \leq b_i,\\
		0 \leq b_i-a_i^T(Uv_k+\delta).
	\end{gather*}
	Since $j \in I_-$, we have $a_j^T\eta < 0$. Thus,
	\begin{gather*}
		0 \geq (b_i-a_i^T(Uv_k+\delta))(a_j^T\eta).
	\end{gather*}
	Since we have $a_i^T\eta = 0$, $(b_j-a_j^T(Uv_k+\delta))(a_i^T\eta) = 0$. So,
	\begin{gather*}
		0 \geq (b_i-a_i^T(Uv_k+\delta))(a_j^T\eta),\\
		(b_j-a_j^T(Uv_k+\delta))(a_i^T\eta) \geq (b_i-a_i^T(Uv_k+\delta))(a_j^T\eta).\\
	\end{gather*}
	as desired.\\\\
	Now consider the case where $a_i^T\eta > 0$. From the assumptions of this direction, we have
	\begin{gather*}
		a_i^T(Uv_k+\delta) + (a_i^T\eta)z_k \leq b_i,\\
		a_j^T(Uv_k+\delta) + (a_j^T\eta)z_k \leq b_j.
	\end{gather*}
	Since $i \in I_+$ and $j \in I_-$
	\begin{gather*}
		z_k \leq \frac{b_i-a_i^T(Uv_k+\delta)}{a_i^T\eta},\\
		z_k \geq \frac{b_j-a_j^T(Uv_k+\delta)}{a_j^T\eta}.
	\end{gather*}
	Combining these yields
	\begin{gather*}
		\frac{b_j-a_j^T(Uv_k+\delta)}{a_j^T\eta} \leq \frac{b_i-a_i^T(Uv_k+\delta)}{a_i^T\eta},\\
		(b_j-a_j^T(Uv_k+\delta))(a_i^T\eta) \geq (b_i-a_i^T(Uv_k+\delta))(a_j^T\eta).
	\end{gather*}
	As desired.\\\\
	$(\Longleftarrow)$
	\\\\
	By the assumption, we have a bipartition $(I_-,I_+)$ such that for all $1 \leq k \leq n, i \in I_+, j \in I_-,$
	\begin{gather*}
		(b_j-a_j^T(Uv_k+\delta))(a_i^T\eta) \geq (b_i-a_i^T(Uv_k+\delta))(a_j^T\eta).
	\end{gather*}
    and for all $i \in I_+$ and $j \in I_-$, $a_i^T\eta \geq 0$ and $a_j^T \eta \leq 0$.\\\\
	Let $k$ be arbitrary. We want to prove that there exists a $z_k$ such that for all $1 \leq i \leq n$
	\begin{gather*}
		a_i^T(Uv_k+\delta)+z_ka_i^T\eta \leq b.
	\end{gather*}
    For each $1 \leq i \leq n$, we examine the condition above to see the constraints on $z_k$. We then show that, given the assumption, there must exist a solution.\\\\
	Let $1 \leq i \leq n$ be arbitrary. We consider three cases, $a_i^T\eta = 0, a_i^T\eta < 0$ and $a_i^T\eta  > 0$.\\\\
	If $a_i^T \eta = 0$, we then have two more cases, $i \in I_+$ and $i \in I_-$. If $i \in I_+$, the condition becomes
	\begin{gather*}
		a_i^T(Uv_k+\delta)+z_k(0) = a_i^T(Uv_k+\delta) \leq b.
	\end{gather*}
    This condition is independent of $z_k$. So, as long as it is always satisfied, it does not influence the value of $z_k$. Let $j \in I_-$ be such that $a_j^T \eta < 0$. We are guaranteed that one such $j$ must exist since $Q$ is bounded. Using the fact that $a_i^T\eta = 0$, the assumption can be simplified.
    \begin{gather*}
		(b_j-a_j^T(Uv_k+\delta))(a_i^T\eta) \geq (b_i-a_i^T(Uv_k+\delta))(a_j^T\eta),\\
        0 \geq (b_i-a_i^T(Uv_k+\delta))(a_j^T\eta).
	\end{gather*}
    Since $a_j^T\eta < 0$,
    \begin{gather*}
        0 \leq b_i-a_i^T(Uv_k+\delta),\\
        a_i^T(Uv_k+\delta) \leq b_i.
	\end{gather*}
	as a result, our assumptions guarantee this condition is always satisfied. The case where $i \in I_-$ is very similar, but with the roles of $i$ and $j$ in the previous proof reversed.\\\\
	If $a_i^T\eta > 0$, the condition becomes
	\begin{gather*}
		z_k \leq \frac{b_i-a_i^T(Uv_k+\delta)}{a_i^T\eta}.
	\end{gather*}
	If $a_i^T\eta < 0$, the condition becomes
	\begin{gather*}
		z_k \geq \frac{b_i-a_i^T(Uv_k+\delta)}{a_i^T\eta}.
	\end{gather*}
	As a result, such a $z_k$ exists if and only if for all $i \in I_+$ and $j \in I_-$
	\begin{gather*}
		\frac{b_j-a_j^T(Uv_k+\delta)}{a_j^T\eta} \leq \frac{b_i-a_i^T(Uv_k+\delta)}{a_i^T\eta},\\
		(b_j-a_j^T(Uv_k+\delta))(a_i^T\eta) \geq (b_i-a_i^T(Uv_k+\delta))(a_j^T\eta).
	\end{gather*}
    Which holds by the hypothesis.
\end{proof}

\end{document}
